\paragraph{Degeneration Analysis}

For the degeneration analysis, we systematically disrupted the consensus network by progressively rewiring its edges and measured how each distance metric responded to increasing structural change. At each step $k \in \{1,\ldots,K\}$ with $K=100$, a rewired network was generated independently from the original consensus by attempting $r_k = \lfloor M \cdot k / K \rfloor$ double-edge swaps (using NetworkX's \texttt{double\_edge\_swap} function), where $M$ is the total number of edges. This preserves the degree sequence of every node while randomising the global wiring pattern. Importantly, each rewired network is generated fresh from the original consensus — not cumulatively from the previous step — so that the rewiring fraction increases monotonically from nearly $0$ at $k=1$ to $1$ at $k=K$, representing the full spectrum from the original consensus to a maximally randomised network.

Because \texttt{double\_edge\_swap} attempts swaps that can occasionally fail or cancel, the number of attempted swaps $r_k$ does not exactly equal the number of edges that actually changed. We therefore quantified the effective degeneration at each step as the Hamming distance $h_k = |A - A_k'|$ between the original adjacency matrix $A$ and the rewired adjacency matrix $A_k'$, which directly counts the number of edges that differ between the two networks. To account for stochasticity in the rewiring procedure, we ran 200 independent processes per metric, each starting from the original consensus network but using a different random seed.

To assess how faithfully each distance metric tracks effective structural disruption, we computed the mean absolute error (MAE) between each metric and Hamming distance, which served as a reference measure of true edge-level change. Prior to computing the MAE, both the metric values and the Hamming distances were min-max normalised globally to $[0,1]$ across all observations (all steps and all processes combined). Metrics defined as similarities rather than distances were additionally inverted ($1 - \text{value}$) to convert them into distances before normalisation. The MAE for each metric was then computed as:
$$\mathrm{MAE}=\frac{1}{N}\sum_{n=1}^N |y_n - h_n|,$$
where $y_n$ and $h_n$ denote the normalised metric value and the normalised Hamming distance for the $n$-th process-step observation, respectively, and $N$ is the total number of observations across all processes and steps. A lower MAE indicates that the metric's response to rewiring closely mirrors the actual number of changed edges, while a higher MAE reflects a non-linear or otherwise distorted sensitivity profile.
